\section{Introduction}
\label{sec:intro}

Distribution networks with high penetrations of Distributed Energy
Resources (DERs) --- rooftop photovoltaics, battery storage, controllable
loads --- are part of critical national energy infrastructure, and the
cyber surface that orchestrates them is broad: DER management systems
(DERMS) issue dispatch commands, inverters and intelligent electronic
devices report telemetry over IEC~61850, DNP3, Modbus/TCP and SunSpec, and
aggregator platforms federate fleets across utility
boundaries~\cite{CaballeroPena2022_DER_DN,Zografopoulos2023_DER_CybersecurityOutlook}.
The U.S. Department of Energy projects that DER growth will make this
surface a grid-operations problem within the
decade~\cite{DOE_CESER_DER2022}, and current standards work --- IEEE
1547.3-2023 for DER cybersecurity~\cite{IEEE15473_2023}, NERC CIP-015-1
for monitoring inside the OT perimeter~\cite{NERC_CIP015_2024}, NIST SP
800-82r3 and SP 800-61r3 for OT security and incident
response~\cite{NIST80082r3,NIST80061r3} --- consistently stops at
\emph{guidance}: none of it prescribes the run-time mechanism that turns a
detected DER incident into a safe physical mitigation.

That gap is where large language models are now being proposed. LLM-based
SOC automation promises to close the triage-to-response distance, but it
imports failure modes that operational technology cannot absorb by
default: hallucinated commands, prompt injection through
attacker-controlled log lines, non-deterministic reasoning, and
schema-violating outputs~\cite{Yao2024_LLM_SecPrivacy,MITRE_ATLAS}. In an
IT SOC these are quality problems; in a feeder control path they are
physical hazards, because a wrong action --- isolating a healthy inverter,
say --- sheds real generation. We therefore ask a narrower question than
whether an LLM \emph{improves} DER incident response:

\begin{quote}
\emph{Can an unverified LLM be consulted safely at all inside a
cyber-physical incident-response path --- and which mechanisms actually
carry the safety when it is?}
\end{quote}

This paper answers with a runtime-assurance architecture,
\ours{}, in the
lineage of the Simplex construction~\cite{Sha2001Simplex} and shielded
decision-making~\cite{Alshiekh2018Shielding}: the language model is an
\emph{untrusted advisory component}, and a small deterministic layer owns
the executable action surface. Concretely, a deterministic FeatureView
summarises telemetry and protocol evidence; an \emph{Incident Evidence
Gate} authorizes response --- and decides whether the model is consulted at
all --- when prespecified observable protocol evidence satisfies its frozen
persistence thresholds; a QLoRA-adapted
open-weights model (Qwen2.5-7B or Llama-3.1-8B) proposes a hypothesis and
an action; and a runtime-assurance shield --- a fixed five-primitive
registry, class-aware removals, a rule that a model-proposed irreversible
action never auto-executes (it is vetoed or escalated to a human;
Table~\ref{tab:action-policy}), and a counterfactually validated
horizon-aware impact estimator that ranks the admissible candidate set on
fallback ---
bounds what can execute. A deterministic fast path answers immediately in
parallel; the model path runs under an explicit latency budget.

Everything we claim is measured with the real model in the loop, under
serving assertions that abort any run in which a fallback, a wrong
adapter, or missing provenance would otherwise masquerade as an LLM
result, and with every runtime feature computed only from wire-observable
telemetry --- a provenance discipline we audit feature by feature
(\S\ref{sec:method-featureview}). The evaluation is
organised around four research questions: whether adversarial or incorrect
model output reaches an unsafe executable action (RQ1); whether the
Evidence Gate reduces unnecessary benign response while retaining
independently evidenced incidents (RQ2); what physical cost remains after
containment on a held-out set of configurations (RQ3); and whether
the safety boundary holds under realistic latency, fallback, and
escalation (RQ4). It spans a canonical 70-case leakage-controlled
adversarial suite, a runtime-safe gate robustness evaluation, a
hash-frozen held-out 49-configuration end-to-end set, and
a counterfactual estimator validation --- all with the configuration, not
the seed, as the unit of inference.

The findings are deliberately asymmetric. The evaluated LLM advisers do not
improve physical outcomes: with the shield in place, the projection arms show no
statistically detectable difference from the deterministic fast path (holdout
mean ENS $0.54$/$0.75$\,kWh for Qwen/Llama versus $0.35$; both Holm $p=1.0$), and
a ground-truth-class oracle yields only a \emph{modest} improvement
($-0.29$\,kWh, 95\% CI $[-0.69,-0.05]$, not surviving multiplicity correction)
that the LLMs do not reproduce. What the architecture buys is \emph{containment}:
against the adversarial suite the model is consulted on the majority of cases yet
the shielded architecture executes an unsafe irreversible action zero times,
against $61/70$ (unshielded Qwen) and $18/70$ (unshielded Llama); an ablation
localises this in the class-aware avoidance, and an exhaustive enumeration of
$40{,}824$ modelled decision states finds no invariant violation while
guard-removal mutations do. We also report results we did not want: a
within-capacity false-data-injection attack has no observable hard signature and
is not separable from benign operation without redundant sensing; and two shield
designs were falsified by our own evaluation --- gating an irreversible action on
model-stated confidence was \emph{worse than no shield}, and honouring a model
\textsc{no\_op} let a denial-of-service outage run unmitigated --- together, the
reason model output must gate neither action nor inaction.
\ours{} does not establish that an LLM improves DER incident response; it
establishes an execution boundary under which an LLM adviser can be evaluated
without granting it direct physical authority.

\paragraph{Contributions}
\begin{enumerate}
  \item \textbf{A deterministic runtime-assurance architecture, and its
    safety invariants, for bounded LLM-assisted DER incident response}
    (\S\ref{sec:architecture}). Evidence-threshold response authorization, action
    bounding, irreversible-action escalation, and the symmetric rule that model
    output may gate neither action nor inaction are deterministic and each
    stated as a directly testable invariant; every evidence feature is
    computed from wire-observable telemetry and audited against a provenance
    table that records each feature's source and the attacker's influence
    over it.
  \item \textbf{A proposal-versus-execution adversarial evaluation plus
    exhaustive structural verification} (\S\ref{sec:rq1}, \S\ref{sec:rq4}). A
    leakage-controlled 70-case suite with per-decision provenance separates
    proposal from execution; an exhaustive enumeration of $40{,}824$ modelled
    decision states with guard-removal mutation tests shows the safety
    invariants are load-bearing; and two falsified designs ---
    confidence-gated autonomy and honoured stand-down --- show why
    attacker-influenceable model output must not relax constraints in either
    direction.
  \item \textbf{Independent cyber-physical, gate, estimator, and timing
    validation with configuration-level inference}
    (\S\ref{sec:rq2}--\S\ref{sec:rq4}). A hash-frozen held-out 49-configuration end-to-end set of
    configurations drawn from prespecified axes, a runtime-safe gate robustness
    evaluation reporting the honest per-family detectability the observable-only
    feature set exposes, a counterfactually validated estimator
    stress-tested under horizon and execution perturbations, and a
    deadline-aware latency analysis.
\end{enumerate}

Section~\ref{sec:background} states the setting, threat model, and trust
boundary; Section~\ref{sec:relatedwork} positions the work;
Section~\ref{sec:architecture} describes \ours{} and its invariants;
Section~\ref{sec:eval} reports the experiments around four research
questions; Section~\ref{sec:discussion} discusses what does and does not
provide safety; Section~\ref{sec:conclusion} concludes.
